\frontmatterchapter{SUMMARY}

This thesis addresses the problem of autonomous interaction between a mobile manipulator and its
environment, using the opening of two types of door, revolute and sliding, as the test case. A
complete system integrating perception, force and torque estimation, and contact control was
developed in simulation for the KUKA KMR iiwa mobile manipulator. The perception module
determines the position and orientation of the door handle from camera images using visual
markers, while the force and torque at the tool tip are estimated from joint torques, without a
dedicated force sensor. The door type is not supplied to the robot as a parameter, but inferred
from the geometry of the handle.

The control problem is solved by two approaches. The classical, model-based approach performs the
entire task chain, from approach and grasp through opening to passing through the doorway,
relying on the constraint geometry computed from perception. The learning-based approach
formulates the task as a Markov decision process with variable impedance as the action space, so
that the policy selects the controller stiffness along each axis at every step, and must infer
the constraint geometry from its own force and displacement observations.

Both approaches were evaluated on the same robot, the same scene and the same success criteria.
For sliding doors, both reach the full target in every trial. For revolute doors, the classical
approach reliably opens the door to 50\textdegree{} but cannot exceed approximately
57\textdegree{}, as the handle slips out of the gripper, whereas the learned policy opens the
door beyond 90\textdegree{} in 96 \% of episodes, with substantially lower contact forces. The
conclusion is that explicit knowledge of the constraint geometry is sufficient when that geometry
is simple and well estimated, but loses its advantage as soon as the task requires the stiffness
to be adapted during the interaction itself.

\vspace{8mm}
\noindent\textbf{Key words:} mobile manipulator, physical interaction, environment property
estimation, impedance control, reinforcement learning, simulation.
\clearpage